\documentclass[11pt,letterpaper]{article}

% Modern fonts and typography
\usepackage{lmodern}
\usepackage{tgpagella}
\usepackage[T1]{fontenc}
\usepackage[utf8]{inputenc}
\usepackage{microtype}

% Layout and spacing
\usepackage[margin=1in,top=0.75in,bottom=0.75in]{geometry}
\usepackage{parskip}
\usepackage{setspace}

% Graphics and colors
\usepackage{graphicx}
\usepackage[dvipsnames,table]{xcolor}
\usepackage{tikz}
\usetikzlibrary{shapes,decorations,shadows}

% Tables and lists
\usepackage{tabularx}
\usepackage{booktabs}
\usepackage{longtable}
\usepackage{enumitem}
\usepackage{multicol}
\usepackage{pbox}

% Headers and footers
\usepackage{fancyhdr}
\usepackage{lastpage}

% Sections and formatting
\usepackage{titlesec}
\usepackage{fontawesome5}

% Hyperlinks
\usepackage{hyperref}

% Colored boxes
\usepackage{tcolorbox}
\tcbuselibrary{skins,breakable}

% Define modern color scheme
\definecolor{primaryblue}{RGB}{0,47,108}
\definecolor{accentblue}{RGB}{44,130,201}
\definecolor{darkgray}{RGB}{64,64,64}
\definecolor{lightgray}{RGB}{245,245,245}
\definecolor{urgold}{RGB}{255,204,0}
\definecolor{nysedbg}{RGB}{240,248,255}

% Hyperlink setup
\hypersetup{
    colorlinks=true,
    linkcolor=accentblue,
    urlcolor=accentblue,
    citecolor=primaryblue,
    pdftitle={LING 214/414: Statistical Methods in Linguistics},
    pdfauthor={Aaron Steven White}
}

% Custom font settings
\renewcommand{\familydefault}{\sfdefault}
\setstretch{1.05}

% Section formatting with modern style
\titleformat{\section}
    {\Large\bfseries\color{primaryblue}}
    {\thesection}
    {1em}
    {}
    [\vspace{0.3em}\titlerule]

\titleformat{\subsection}
    {\large\bfseries\color{darkgray}}
    {\thesubsection}
    {1em}
    {}

% Custom commands
% NYSED tag removed

\newcommand{\courseinfo}[2]{%
    \noindent\makebox[0.3\textwidth][l]{\textbf{#1}} #2
}

% Header and footer setup
\pagestyle{fancy}
\fancyhf{}
\fancyhead[L]{\small\color{darkgray}LING 214/414-01 | Fall 2026}
\fancyhead[R]{\small\color{darkgray}Statistical Methods in Linguistics}
\fancyfoot[C]{\small\color{darkgray}Page \thepage\ of \pageref{LastPage}}
\renewcommand{\headrulewidth}{0.5pt}
\renewcommand{\footrulewidth}{0pt}

\begin{document}

% Modern title section
\begin{center}
    \colorbox{primaryblue}{\parbox{\textwidth}{%
    \vspace{0.5cm}
    \begin{center}
    \color{white}
    {\Huge\bfseries Statistical Methods in Linguistics}\\[0.3em]
    {\Large LING 214 / LING 414 \quad\textbullet\quad Fall 2026}\\[0.5em]
    \rule{0.5\textwidth}{0.5pt}\\[0.3em]
    {\large University of Rochester}\\
    {\large Department of Linguistics}
    \end{center}
    \vspace{0.5cm}
    }}
\end{center}

\vspace{1cm}

% Instructor information box
\begin{tcolorbox}[
    colback=lightgray,
    colframe=darkgray,
    boxrule=1pt,
    arc=2mm,
    left=10pt,
    right=10pt,
    top=10pt,
    bottom=10pt
]
\begin{minipage}[t]{0.48\textwidth}
    \textbf{\large\color{primaryblue}Instructor}\\[0.3em]
    \faUser\ \textbf{Aaron Steven White}\\
    \faMapMarker*\ 511A Lattimore Hall\\
    \faGlobe\ \href{http://aaronstevenwhite.io}{aaronstevenwhite.io}\\
    \faCalendarCheck\ \href{https://app.squarespacescheduling.com/schedule.php?owner=26977798&appointmentType=36264617}{Office Hours Scheduler}
\end{minipage}
\hfill
\begin{minipage}[t]{0.48\textwidth}
    \textbf{\large\color{primaryblue}Course Meetings}\\[0.3em]
    \faClock\ Monday \& Wednesday\\
    \phantom{\faClock\ }12:30 PM to 1:45 PM\\
    \faMapMarker*\ Lattimore 513\\
    \faComments\ \href{https://stats-in-ling-fall-2026.zulipchat.com/}{Course Zulip}\\
\end{minipage}
\end{tcolorbox}

\section{Course Overview}

\subsection{Course Communication}

\begin{tcolorbox}[
    colback=urgold!10,
    colframe=urgold,
    boxrule=1pt,
    arc=2mm,
    left=10pt,
    right=10pt
]
\faExclamationTriangle\ \textbf{Important: All course communication will be conducted through Zulip}\\

\faComments\ \textbf{Course Zulip:} \href{https://stats-in-ling-fall-2026.zulipchat.com/}{stats-in-ling-fall-2026.zulipchat.com}
\begin{itemize}[leftmargin=*,itemsep=2pt]
    \item Post general questions to the appropriate channel
    \item Use direct messages (DMs) for private matters
    \item Do not use email for course-related communication
\end{itemize}
\vspace{1em}
If you are registered for the course but cannot access the workspace, please contact me.
\end{tcolorbox}

\subsection{Course Description}

This course provides an introduction to probability and statistics for linguistics, serving as an essential foundation for linguistics students who aim to analyze experimental and corpus linguistic data. The course covers:

\begin{multicols}{2}
\begin{itemize}[leftmargin=*,itemsep=0pt]
    \item Elementary probability theory
    \item Descriptive and inferential statistics
    \item Fixed and mixed effects models
    \item Model evaluation and comparison
    \item Analysis of corpus data
    \item Analysis of judgment data
    \item Analysis of reading time data
    \item Analysis of acoustic measures
\end{itemize}
\end{multicols}

\subsection{Prerequisites}

This course is restricted to students who have completed \textbf{LING 110} with a grade of C- or better.

\subsection{Credit Hours}

This 4-credit course includes 150 minutes per week of direct instruction (two 75-minute sessions) and a minimum of 480 minutes per week of out-of-class student work, including reading assignments, problem sets, data analysis, and project work. This aligns with the University of Rochester's \href{https://www.rochester.edu/bulletin/policies/academic/credit-hour/}{credit hour policy}, where each credit hour requires 50 minutes of instruction and 120 minutes of supplementary work.

As a cross-listed course, LING414 students are expected to demonstrate advanced content mastery, rigor, and requirements beyond the LING214 level. Graduate students will complete a comprehensive final project that demonstrates independent research capabilities consistent with graduate-level scholarship.

\section{Learning Objectives and Outcomes}

\subsection{Course Objectives}

By the end of this course, students will:

\begin{enumerate}[leftmargin=*]
    \item \textbf{Master} fundamental concepts in probability and statistics
    \item \textbf{Apply} statistical methods to linguistic data analysis
    \item \textbf{Design} and analyze linguistic experiments
    \item \textbf{Interpret} statistical results in linguistics research
    \item \textbf{Implement} analyses using R programming
    \item \textbf{Construct} mixed effects models for complex datasets
\end{enumerate}

\section{Course Materials}

\subsection{Textbooks}

\begin{tcolorbox}[colback=white,colframe=accentblue,boxrule=0.5pt,arc=2mm]
\faBook\ \textbf{An Introduction to Bayesian Data Analysis for Cognitive Science}\\
\phantom{\faBook\ }Bruno Nicenboim, Daniel Schad, and Shravan Vasishth\\
\phantom{\faBook\ }\href{https://bruno.nicenboim.me/bayescogsci/}{bruno.nicenboim.me/bayescogsci}
\end{tcolorbox}

\begin{tcolorbox}[colback=white,colframe=accentblue,boxrule=0.5pt,arc=2mm]
\faBook\ \textbf{Statistics for Linguists: An Introduction Using R}\\
\phantom{\faBook\ }Bodo Winter\\
\phantom{\faBook\ }\href{https://appliedstatisticsforlinguists.org/bwinter_stats_proofs.pdf}{appliedstatisticsforlinguists.org/bwinter\_stats\_proofs.pdf}
\end{tcolorbox}

\begin{tcolorbox}[colback=white,colframe=accentblue,boxrule=0.5pt,arc=2mm]
\faBook\ \textbf{Doing Bayesian Data Analysis: A Tutorial with R, JAGS, and Stan}\\
\phantom{\faBook\ }John K. Kruschke, second edition (2015)\\
\phantom{\faBook\ }\href{https://www.educate.elsevier.com/book/details/9780124058880}{Elsevier book page}
\end{tcolorbox}

\subsection{Software}

\begin{itemize}[leftmargin=*]
    \item \faCode\ \textbf{R} | Statistical computing language\\
    \phantom{\faLaptop\ }\href{https://www.r-project.org/}{www.r-project.org}
    
    \item \faLaptop\ \textbf{RStudio} | Integrated development environment\\
    \phantom{\faCode\ }\href{https://posit.co/download/rstudio-desktop/}{posit.co/download/rstudio-desktop}

    \item \faTerminal\ \textbf{Visual Studio Code} | Integrated development environment\\
    \phantom{\faTerminal\ }\href{https://code.visualstudio.com/}{code.visualstudio.com/}
\end{itemize}

\subsection{Readings}

The readings are introduced in the chapter where they become useful. Each chapter contains a callout that explains what to take from the reading and how it bears on the analysis developed there.

\begin{tcolorbox}[
    colback=white,
    colframe=accentblue,
    boxrule=0.5pt,
    arc=2mm,
    breakable,
    left=8pt,right=8pt,top=6pt,bottom=6pt
]
\small
\begin{itemize}[leftmargin=*,itemsep=5pt,topsep=0pt]
    \item \href{https://aclanthology.org/2020.lrec-1.699/}{\textbf{The Universal Decompositional Semantics Dataset and Decomp Toolkit}}. White et al. (2020), Section 2. \textit{Linguistic data.}
    \item \href{https://doi.org/10.3758/s13423-014-0585-6}{\textbf{Zipf's word frequency law in natural language: A critical review and future directions}}. Piantadosi (2014). \textit{Zipf's law.}
    \item \href{https://doi.org/10.1017/CBO9781139013734.006}{\textbf{Population samples}}. Buchstaller \& Khattab (2013), Chapter 5 of \textit{Research Methods in Linguistics}. \textit{Populations and samples.}
    \item \href{https://doi.org/10.1121/1.411872}{\textbf{Acoustic characteristics of American English vowels}}. Hillenbrand, Getty, Clark, \& Wheeler (1995), speaker sample, recording procedure, and formant measurements. \textit{Problem Set 1.}
    \item \href{https://bruno.nicenboim.me/bayescogsci/ch-cv.html}{\textbf{Cross-validation}}. Nicenboim, Schad, \& Vasishth, chapter opening and Section 14.2 through Equation 14.6 of \textit{An Introduction to Bayesian Data Analysis for Cognitive Science}. \textit{Cross-validation.}
    \item \href{https://aclanthology.org/L18-1012/}{\textbf{The Natural Stories Corpus}}. Futrell et al. (2018), Sections 3.1 and 3.3. \textit{Grouped cross-validation.}
    \item \href{https://doi.org/10.1016/j.cogpsych.2016.06.002}{\textbf{Limits on lexical prediction during reading}}. Luke \& Christianson (2016), read with \href{https://doi.org/10.3758/s13428-017-0908-4}{\textbf{The Provo Corpus}}. Luke \& Christianson (2018), data description. \textit{Problem Set 2.}
    \item \href{https://appliedstatisticsforlinguists.org/bwinter_stats_proofs.pdf}{\textbf{Analyzing linguistic diversity using Poisson regression; adding exposure variables}}. Winter (2020), Sections 13.3 and 13.4 of \textit{Statistics for Linguists}. \textit{Poisson regression and exposure offsets.}
    \item \href{https://doi.org/10.1007/s11525-024-09429-8}{\textbf{French liaison is allomorphy, not allophony: Evidence from lexical statistics}}. Storme (2024), Sections 4.1 and 4.2.2. \textit{Negative-binomial regression.}
    \item \href{https://doi.org/10.1007/s11050-025-09244-9}{\textbf{Factivity, presupposition projection, and the role of discrete knowledge in gradient inference judgments}}. Grove \& White (2025), supplementary material, Section 1. \textit{Why ordinary beta regression fails for sliders.}
    \item \href{https://appliedstatisticsforlinguists.org/bwinter_stats_proofs.pdf}{\textbf{The independence assumption; dealing with non-independence via experimental design and averaging}}. Winter (2020), Sections 14.2 and 14.3 of \textit{Statistics for Linguists}. \textit{Structured linguistic data.}
    \item \href{https://web.stanford.edu/~bresnan/qs-submit.pdf}{\textbf{Predicting the Dative Alternation}}. Bresnan, Cueni, Nikitina, \& Baayen (2007), pages 69 to 94. \textit{Problem Set 3.}
    \item \href{https://ojs.ub.uni-konstanz.de/sub/index.php/sub/article/view/601}{\textbf{The CommitmentBank: Investigating projection in naturally occurring discourse}}. de Marneffe, Simons, \& Tonhauser (2019), response scale, embedding environments, and sources of variation. \textit{Ordinal mixed models and Problem Set 4.}
    \item \href{https://doi.org/10.16995/labphon.10855}{\textbf{Final devoicing before it happens: A large-scale study of word-final obstruents in French}}. Jatteau et al. (2025), selected paragraphs of Section 2.2. \textit{Support check for ordered-beta regression.}
    \item \href{../readings/white-grimm-glass-2026-aspectual-similarity.pdf}{\textbf{Aspectual similarity predicts sense similarity}}. White, Grimm, \& Glass (2026), Experiment 2 through Analysis. \textit{Bounded similarity judgments.}
    \item \href{https://doi.org/10.1121/1.1906875}{\textbf{Control methods used in a study of the vowels}}. Peterson \& Barney (1952), speaker sample, vowel materials, acoustic measurements, and listener task. \textit{Gaussian mixtures.}
    \item \href{https://doi.org/10.5334/gjgl.1001}{\textbf{Frequency, acceptability, and selection: A case study of clause-embedding}}. White \& Rawlins (2020), Sections 1, 2.4, 5, and 6 and Appendix C. \textit{Distributional prediction and factorization.}
    \item \href{https://aclanthology.org/2022.tacl-1.2/}{\textbf{Decomposing and recomposing event structure}}. Gantt, Glass, \& White (2022), data and model sections. \textit{Custom annotation models.}
\end{itemize}
\end{tcolorbox}

\section{Course Structure and Requirements}

\subsection{Assessment Overview}

\begin{center}
\begin{tcolorbox}[
    colback=white,
    colframe=accentblue,
    boxrule=0.5pt,
    arc=2mm,
    width=.92\textwidth
]
\centering
\renewcommand{\arraystretch}{1.4}
\begin{tabular}{l>{\centering}p{1.8cm}>{\centering}p{1.8cm}p{6cm}}
\rowcolor{accentblue!10}
\textbf{Component} & \textbf{LING 214} & \textbf{LING 414} & \textbf{Details}\\
\midrule
\faEdit\ Problem Sets & \textbf{65\%} & \textbf{40\%} & Five coding assignments with in-class assessments\\
\faGraduationCap\ Midterm Exam & \textbf{25\%} & \textbf{20\%} & Comprehensive exam on Oct. 28\\
\faUsers\ Participation & \textbf{10\%} & N/A & Attend office hours twice\\
\faTasks\ Final Project & N/A & \textbf{40\%} & See rubric\\
\bottomrule
\end{tabular}
\end{tcolorbox}
\end{center}

\subsection{Problem Set Grading}

\begin{tcolorbox}[
    colback=lightgray,
    colframe=darkgray,
    boxrule=1pt,
    arc=2mm,
    left=10pt,
    right=10pt
]
\faEdit\ \textbf{Problem Set Assessment Structure}\\[0.3em]
Each problem set will be graded on two components:
\vspace{0.2em}
\begin{itemize}[leftmargin=*,itemsep=2pt]
    \item \textbf{50\%: Functionality and Completeness:} Code works correctly and addresses all requirements
    \item \textbf{50\%: In-Class Presentation:} Randomly selected students present their homework solution and explain how it works to the class. All students must be prepared to present on any assessment day following the problem set due date.
\end{itemize}
\vspace{0.5em}
\faExclamationTriangle\ \textbf{Important Notes:}
\begin{itemize}[leftmargin=*,itemsep=2pt]
    \item Students will not know in advance if they will present on a given day
    \item All students must attend class and be ready to present their solutions
    \item \textbf{LLM Usage:} While LLMs may be used for problem sets, you must fully understand your code to present it effectively
\end{itemize}
\end{tcolorbox}

\subsection{Final Project Requirements (LING 414 Only)}

\begin{tcolorbox}[
    colback=urgold!10,
    colframe=urgold,
    boxrule=1pt,
    arc=2mm
]
\faGraduationCap\ \textbf{Final Project Components (40\% of total grade)}\\[-1em]
\begin{itemize}[leftmargin=*,itemsep=2pt]
    \item \textbf{Pre-proposal meeting (5\%):} Meeting with me to discuss proposal (by Oct. 16)
    \item \textbf{Proposal (5\%):} Proposal outlining research question and methods (due Oct. 28)
    \item \textbf{Presentation (20\%):} In-class presentation with professional slides (Dec. 7, 9, or 14)
    \item \textbf{Code Submission (20\%):} All code structured following the rubric (due Dec. 14); evaluated via a 10 to 20 minute \faComments\ oral assessment scheduled during finals week
    \item \textbf{Paper Submission (50\%):} Research paper (due Dec. 14)
\end{itemize}
\end{tcolorbox}
\begin{tcolorbox}[
    colback=urgold!10,
    colframe=urgold,
    boxrule=1pt,
    arc=2mm
]
\textbf{Project Options:}
\begin{enumerate}[leftmargin=*,itemsep=2pt]
    \item \textbf{Corpus Analysis:} using open source or LDC corpora with external NLP tools
    \item \textbf{Experiment Design:} judgment/reading experiment, including data collection instrument, with simulated data analysis
    \item \textbf{Mechanistic Interpretation:} deep learning model analysis
\end{enumerate}
\vspace{1em}
\faBook\ \textit{See separate Final Project Rubric document for detailed requirements and grading criteria}
\end{tcolorbox}

\subsection{Grading Scale}

\begin{center}
\renewcommand{\arraystretch}{1.3}
\begin{tabular}{*{11}{>{\centering\arraybackslash}p{1.1cm}}}
\toprule
\cellcolor{accentblue!30}\textbf{A} & 
\cellcolor{accentblue!25}\textbf{A-} & 
\cellcolor{accentblue!20}\textbf{B+} & 
\cellcolor{accentblue!15}\textbf{B} & 
\cellcolor{accentblue!10}\textbf{B-} & 
\cellcolor{lightgray}\textbf{C+} & 
\cellcolor{lightgray}\textbf{C} & 
\cellcolor{lightgray}\textbf{C-} & 
\cellcolor{lightgray}\textbf{D+} & 
\cellcolor{lightgray}\textbf{D} & 
\cellcolor{lightgray}\textbf{E}\\
\midrule
93+ & 90 & 87 & 83 & 80 & 77 & 73 & 70 & 67 & 60 & <60\\
\bottomrule
\end{tabular}
\end{center}

\section{Course Policies}

\subsection{Office Hours}

\faCalendarCheck\ \textbf{By appointment} Please schedule using this form:
\href{https://app.squarespacescheduling.com/schedule.php?owner=26977798&appointmentType=36264617}{app.squarespacescheduling.com}

I am available to meet if and only if a time is listed as available on the scheduling form.

\subsection{Communication}

\textbf{All communication must be done through the \href{https://stats-in-ling-fall-2026.zulipchat.com/}{Fall 2026 course Zulip}.}

\begin{itemize}[leftmargin=*]
    \item \faComments\ Post general questions to appropriate Zulip channels
    \item \faUserLock\ Use Zulip DMs for grade questions, course absence notifications, etc.
    \item \faClock\ Response time: 1-2 business days during work hours
    \item \faExclamationTriangle\ No monitoring after 5 PM or weekends
\end{itemize}

\subsection{Late Work Policy}

\begin{itemize}[leftmargin=*]
    \item Problem sets: 10\% deduction per day (max 3 days)
    \item Project components (414 only): No late submissions without prior approval
    \item Extensions require 48-hour advance notice (except emergencies)
\end{itemize}

\subsection{Academic Integrity}

Collaboration on problem sets with classmates is encouraged, but each student must write their own solutions. See the generative AI policy below and see: \href{http://www.rochester.edu/college/honesty/}{rochester.edu/college/honesty}

\subsection{Generative AI Policy}

Students may use generative AI tools (e.g., ChatGPT, Claude, GitHub Copilot) for coding assistance on problem sets and projects. However, students must:
\begin{itemize}[leftmargin=*]
    \item Fully understand and be able to explain all generated code. This will be assessed during in-class presentations and oral assessments.
    \item Cite in code comments any AI tools used, including the specific prompts
    \item Not use AI tools for written assignments
\end{itemize}

\subsection{Accessibility}

The University of Rochester welcomes students with disabilities. Students seeking accommodations must request them through the Office of Disability Resources: \href{https://www.rochester.edu/disability/}{rochester.edu/disability}

\subsection{Title IX}

All faculty are mandatory reporters for Title IX issues. Students can contact the Title IX Office directly at \href{mailto:titleix@rochester.edu}{titleix@rochester.edu} or 585-275-1654. More information: \href{www.rochester.edu/sexualmisconduct}{rochester.edu/sexualmisconduct}

\section{Course Schedule}

\begingroup
\footnotesize
\setlength{\tabcolsep}{4pt}
\renewcommand{\arraystretch}{1.45}
\begin{longtable}{>{\centering}p{0.8cm}p{2.3cm}p{8.4cm}>{\centering\arraybackslash}p{2.8cm}}
\toprule
\rowcolor{gray!90}
\textcolor{white}{\textbf{Week}} & \textcolor{white}{\textbf{Dates}} & \textcolor{white}{\textbf{Topics}} & \textcolor{white}{\textbf{Due}}\\
\midrule
\endfirsthead
\multicolumn{4}{c}{\textit{...continued from previous page}}\\
\toprule
\rowcolor{gray!90}
\textcolor{white}{\textbf{Week}} & \textcolor{white}{\textbf{Dates}} & \textcolor{white}{\textbf{Topics}} & \textcolor{white}{\textbf{Due}}\\
\midrule
\endhead
\midrule
\multicolumn{4}{r}{\textit{Continued on next page...}}\\
\endfoot
\bottomrule
\endlastfoot

% \faRocket\ 
% \faUsers\ 
% \faUsers\ 

\rowcolor{lightgray!30}
1 & Aug. 31, Sep. 2 & Linguistic data and representation; probability spaces & \\
\rowcolor{red!15}
2 & Sep. 7, 9 & \textit{Labor Day}; events and probability measures & \\
\rowcolor{lightgray!30}
3 & Sep. 14, 16 & Random variables; discrete distributions & \\
4 & Sep. 21, 23 & Continuous and joint distributions; dependence & \\
\rowcolor{lightgray!30}
5 & Sep. 28, 30 & Populations, samples, estimands, and uncertainty & \\
6 & Oct. 5, 7 & Paired inference; conditional means and regression & \textbf{PS1 (Oct. 7)}\\
\rowcolor{red!15}
7 & Oct. 12, 14 & \textit{Fall Break}; multiple regression and model criticism & \\
\rowcolor{lightgray!30}
8 & Oct. 19, 21 & Out-of-sample prediction, loss, and grouped validation & \textbf{PS2 (Oct. 21)}\\
\rowcolor{yellow!20}
9 & Oct. 26, 28 & Review; \textbf{\faEdit\ Midterm} & \textbf{Proposal (414)}\\
\rowcolor{lightgray!30}
10 & Nov. 2, 4 & GLMs: binary, count, and bounded responses & \\
\rowcolor{lightgray!30}
11 & Nov. 9, 11 & Experimental design, dependence, and generalization & \textbf{PS3 (Nov. 11)} \\
12 & Nov. 16, 18 & Mixed models: ordinal and bounded responses & \\
\rowcolor{red!15}
13 & Nov. 23, 25 & Clustering and mixture models; \textit{Thanksgiving recess} & \textbf{PS4 (Nov. 23)}\\
14 & Nov. 30, Dec. 2 & Factorization and missing-data validation; custom model design & \\
\rowcolor{accentblue!15}
15 & Dec. 7, 9 & \faDesktop\ Presentations (414) & \textbf{PS5 (Dec. 7)}\\
\rowcolor{accentblue!15}
16 & Dec. 14 & \faDesktop\ Presentations (414) & \textbf{Paper \& code (414)}\\
\rowcolor{orange!20}
 & Dec. 18 to 23 & \faComments\ Oral assessments (414) & \\
\end{longtable}
\endgroup

\clearpage
\section{Student Resources}

\begin{multicols}{2}
\subsection{Academic Support}
\begin{itemize}[leftmargin=*,itemsep=2pt]
    \item Center for Excellence in Teaching
    \item Writing, Speaking \& Argument Program
    \item Q\&i Statistics Help
    \item Library Research Support
\end{itemize}

\subsection{Health \& Wellness}
\begin{itemize}[leftmargin=*,itemsep=2pt]
    \item University Counseling: 585-275-3113
    \item University Health: 585-275-2662
    \item CARE Network
    \item IT Help Desk: 585-275-2000
\end{itemize}
\end{multicols}

\vfill

\begin{center}
\color{darkgray}
\rule{0.3\textwidth}{0.5pt}\\[0.5em]
\small\textit{This syllabus is subject to change. Updates will be announced in class and on Zulip.}\\
Last updated: \today
\end{center}

\end{document}
